\section{Геометрическая и условная вероятность}

\subsection{Геометрическая вероятность}
Классическое определение вероятности неприменимо, если число возможных исходов бесконечно. Для таких случаев вводится геометрическое определение.

Пусть пространство исходов~$\Omega$ представляет собой некоторую область (отрезок, плоскую фигуру, пространственное тело), имеющую конечную меру $\mu(\Omega)$ (длину, площадь, объем соответственно). Если вероятность попадания случайной точки в любую часть области пропорциональна мере этой части и не зависит от ее расположения и формы, то вероятность события $A \subseteq \Omega$ равна:
\[
\Prob(A) = \frac{\mu(A)}{\mu(\Omega)}.
\]

\begin{example}[Задача о спагетти]
Спагетти длиной~1 разрывается в двух случайных точках $x$ и $y$ ($0 \le x \le y \le 1$). Какова вероятность того, что из трех образовавшихся отрезков длинами $x$, $y-x$ и $1-y$ можно сложить треугольник?

\textbf{Решение:} По неравенству треугольника, длина ни одного из отрезков не должна превышать $0{,}5$:
\[
\begin{cases} x \le 0{,}5, \\ y - x \le 0{,}5 \implies y \le x + 0{,}5, \\ 1 - y \le 0{,}5 \implies y \ge 0{,}5. \end{cases}
\]
Область всех возможных исходов $\Omega = \{(x,y): 0 \le x \le y \le 1\}$ имеет площадь $S_\Omega = 0{,}5$. Область благоприятных исходов~$A$ образует треугольник с площадью $S_A = 0{,}125$. Искомая вероятность:
\[
\Prob(A) = \frac{S_A}{S_\Omega} = \frac{0{,}125}{0{,}5} = \frac{1}{4} = 0{,}25.
\]
\end{example}

\begin{example}[Задача о корнях квадратного уравнения]
Коэффициенты $p$ и $q$ выбираются наудачу из отрезка $[-1, 1]$. Какова вероятность того, что квадратное уравнение $x^2 + px + q = 0$ имеет действительные корни?

\textbf{Решение:} Уравнение имеет действительные корни при $D = p^2 - 4q \ge 0 \implies q \le \frac{p^2}{4}$. 
Область $\Omega = [-1,1] \times [-1,1]$, $S_\Omega = 4$. 
Площадь области $A = \{(p,q) \in [-1,1]^2 : q \le \frac{p^2}{4}\}$ вычисляется через интеграл:
\[
S_A = 2 + \int_{-1}^1 \frac{p^2}{4} \, dp = 2 + \frac{1}{6} = \frac{13}{6} \approx 2{,}167.
\]
Вероятность $\Prob(A) = \frac{13/6}{4} = \frac{13}{24} \approx 0{,}5417$.
\end{example}

\subsection{Условная вероятность и теорема умножения}

\begin{definition}
\textbf{Условной вероятностью} $\Prob(A \mid B)$ события~$A$ при условии, что произошло событие~$B$ (где $\Prob(B) > 0$), называется вероятность наступления~$A$ в суженном пространстве исходов, где достоверным стало событие~$B$.
Определяется она формулой:
\[
\Prob(A \mid B) = \frac{\Prob(A \cap B)}{\Prob(B)}.
\]
\end{definition}

Из этого определения напрямую вытекает важная теорема.

\begin{theorem}[Теорема умножения вероятностей]
Вероятность совместного наступления двух событий~$A$ и~$B$ равна произведению вероятности одного из них на условную вероятность другого:
\[
\Prob(A \cap B) = \Prob(B) \cdot \Prob(A \mid B) = \Prob(A) \cdot \Prob(B \mid A).
\]
\end{theorem}

\begin{proof}
В классической вероятностной модели пусть $n$~--- общее число исходов. $m_B$~--- число исходов, благоприятствующих~$B$. $m_{AB}$~--- число исходов, благоприятствующих совместному наступлению $A$ и $B$ ($A \cap B$).
Тогда, если известно, что событие~$B$ произошло, пространство исходов сужается до $m_B$ вариантов. Из них событию~$A$ благоприятствует $m_{AB}$ исходов.
По классическому определению для суженного пространства:
\[
\Prob(A \mid B) = \frac{m_{AB}}{m_B}.
\]
Перепишем дробь, разделив числитель и знаменатель на общее число исходов~$n$:
\[
\Prob(A \mid B) = \frac{m_{AB} / n}{m_B / n} = \frac{\Prob(A \cap B)}{\Prob(B)}.
\]
Умножив обе части на $\Prob(B)$, получаем $\Prob(A \cap B) = \Prob(B) \cdot \Prob(A \mid B)$. Меняя роли $A$ и $B$, аналогично выводится $\Prob(A \cap B) = \Prob(A) \cdot \Prob(B \mid A)$.
\end{proof}

\subsection{Независимость событий}

\begin{definition}
События~$A$ и~$B$ называются \textbf{независимыми}, если наступление одного не меняет вероятности наступления другого, то есть:
\[
\Prob(A \mid B) = \Prob(A).
\]
\end{definition}

Подставляя это в теорему умножения, получаем строгий алгебраический критерий независимости двух событий:
\[
\Prob(A \cap B) = \Prob(A) \cdot \Prob(B).
\]

\subsection{Парадокс Монти Холла}

Парадокс Монти Холла~--- одна из самых известных задач теории вероятностей, наглядно иллюстрирующая контринтуитивную природу условной вероятности.

\begin{example}[Парадокс Монти Холла]
В телеигре участнику предлагают выбрать одну из трех закрытых дверей ($1, 2, 3$). За одной из дверей находится автомобиль, за двумя другими~--- козлы. 
\begin{enumerate}
\item Игрок выбирает одну из дверей (пусть для определенности дверь~1).
\item Ведущий (Монти Холл), который знает, где именно спрятан автомобиль, открывает одну из двух оставшихся дверей (например, дверь~3), за которой гарантированно оказывается козел.
\item Ведущий предлагает игроку выбор: остаться при своем первоначальном решении (дверь~1) или сменить дверь на оставшуюся закрытую (дверь~2).
\end{enumerate}
Выгодно ли сменить свой выбор?

\textbf{Строгое решение через условную вероятность:} \\
Обозначим через $C_i$ гипотезу о том, что автомобиль находится за дверью~$i$ ($i \in \{1, 2, 3\}$). До начала игры все положения равновероятны:
\[
\Prob(C_1) = \Prob(C_2) = \Prob(C_3) = \frac{1}{3}.
\]
Пусть игрок выбрал дверь~1. Обозначим через $M_3$ событие «ведущий открыл дверь~3». 
Определим условные вероятности действий ведущего при различных положениях автомобиля:
\begin{itemize}
\item Если приз за дверью~1 ($C_1$), у ведущего есть выбор между дверями~2 и~3 (за обеими козлы). Он выбирает наудачу: $\Prob(M_3 \mid C_1) = \frac{1}{2}$.
\item Если приз за дверью~2 ($C_2$), ведущий не может открыть выбранную игроком дверь~1 и не может открыть выигрышную дверь~2. У него нет выбора: $\Prob(M_3 \mid C_2) = 1$.
\item Если приз за дверью~3 ($C_3$), ведущий никогда не откроет дверь с автомобилем: $\Prob(M_3 \mid C_3) = 0$.
\end{itemize}

По формуле полной вероятности найдем безусловную вероятность того, что ведущий откроет именно дверь~3:
\[
\Prob(M_3) = \sum_{i=1}^3 \Prob(C_i)\Prob(M_3 \mid C_i) = \frac{1}{3}\cdot\frac{1}{2} + \frac{1}{3}\cdot 1 + \frac{1}{3}\cdot 0 = \frac{1}{6} + \frac{1}{3} = \frac{1}{2}.
\]

Теперь вычислим апостериорные (условные) вероятности выигрыша:
\begin{align*}
\Prob(C_1 \mid M_3) &= \frac{\Prob(C_1 \cap M_3)}{\Prob(M_3)} = \frac{\frac{1}{3} \cdot \frac{1}{2}}{\frac{1}{2}} = \frac{1}{3} \quad \text{(шанс победы без смены двери)}, \\
\Prob(C_2 \mid M_3) &= \frac{\Prob(C_2 \cap M_3)}{\Prob(M_3)} = \frac{\frac{1}{3} \cdot 1}{\frac{1}{2}} = \frac{2}{3} \quad \text{(шанс победы при смене двери)}.
\end{align*}

Следовательно, смена двери удваивает вероятность выигрыша автомобиля с $\frac{1}{3}$ до $\frac{2}{3}$.
\end{example}

\begin{figure}[htbp]
\centering
\begin{tikzpicture}[
  node distance=1.5cm and 2.5cm,
  box/.style={rectangle, draw=blue!70!black, thick, fill=blue!5, rounded corners=3pt, align=center, inner sep=6pt, font=\small},
  action/.style={rectangle, draw=gray!70!black, fill=gray!10, rounded corners=2pt, align=center, inner sep=4pt, font=\footnotesize},
  win/.style={rectangle, draw=teal!70!black, thick, fill=teal!15, rounded corners=2pt, font=\small\bfseries},
  lose/.style={rectangle, draw=red!70!black, thick, fill=red!15, rounded corners=2pt, font=\small\bfseries},
  arrow/.style={-{Stealth[length=2.5mm]}, thick, gray!60!black}
]
  \node[box] (start) {Игрок выбирает дверь~1};

  \node[action, below left=1.2cm and 2.2cm of start] (c1) {Авто за дверью 1\\($p = 1/3$)};
  \node[action, below=1.2cm of start] (c2) {Авто за дверью 2\\($p = 1/3$)};
  \node[action, below right=1.2cm and 2.2cm of start] (c3) {Авто за дверью 3\\($p = 1/3$)};

  \draw[arrow] (start) -- (c1);
  \draw[arrow] (start) -- (c2);
  \draw[arrow] (start) -- (c3);

  \node[action, below=0.8cm of c1] (host1) {Ведущий открывает\\дверь 2 или 3};
  \node[action, below=0.8cm of c2] (host2) {Ведущий обязан\\открыть дверь 3};
  \node[action, below=0.8cm of c3] (host3) {Ведущий обязан\\открыть дверь 2};

  \draw[arrow] (c1) -- (host1);
  \draw[arrow] (c2) -- (host2);
  \draw[arrow] (c3) -- (host3);

  \node[lose, below=0.8cm of host1] (res1) {Смена $\to$ Проигрыш};
  \node[win, below=0.8cm of host2] (res2) {Смена $\to$ Автомобиль!};
  \node[win, below=0.8cm of host3] (res3) {Смена $\to$ Автомобиль!};

  \draw[arrow] (host1) -- (res1);
  \draw[arrow] (host2) -- (res2);
  \draw[arrow] (host3) -- (res3);
\end{tikzpicture}
\caption{Дерево исходов при стратегии обязательной смены двери в парадоксе Монти Холла}
\label{fig:monty_hall}
\end{figure}

\textbf{Интуитивное объяснение:} В момент первоначального выбора игрок ошибается с вероятностью $\frac{2}{3}$. Если игрок выбрал козла (что происходит в $2$ случаях из $3$), у ведущего нет свободы маневра: он обязан открыть единственную оставшуюся дверь с другим козлом. Тем самым ведущий непреднамеренно указывает игроку пальцем на дверь с автомобилем. Поэтому смена выбора выигрывает ровно тогда, когда исходный выбор был ошибочным, то есть с вероятностью $\frac{2}{3}$.